\section{Related work}
\label{sec:relatedwork}
\subsection{Continual Learning}

There are many continual learning setups that have been discussed recently~\cite{prabhu2020gdumb}. First, continual learning can be divided into task-based and task-free setups. Task-based continual learning~\cite{delange2021continual,zhao2023adaptcl} assumes that task-ID can be available while the network in task-free continual learning~\cite{zhou2021learning} has no access to task-ID at inference time. Next, continual learning can be categorized into offline and online setups. Under the offline setup~\cite{sun2020continual}, all training samples in the current learning stage can be learned with multiple training steps. In contrast, each training sample in the streamed data is seen only once for the online setup~\cite{mai2022online}. Finally, following~\cite{koh2021online}, continual learning also can be separated into disjoint-task and blurry-task setups. The disjoint-task setup requires that the intersection of the label sets of any two learning stages is an empty set~\cite{aljundi2019task}. The blurry-task setup has been proposed to set some common samples for the training samples of each learning stage~\cite{bang2021rainbow}. 


Five main directions drive recent advances in continual learning. 1) {Architecture-based methods} divide each task into a set of specific parameters of the model. They dynamically extend the model as the task increases~\cite{gao2022efficient} or gradually freeze part of parameters to overcome the forgetting problem~\cite{miao2021continual}. 2) {Regularization-based methods} store the historical information learned from old data as the prior knowledge of the network. It consolidates past knowledge by extending the loss function with an additional regularization term~\cite{kirkpatrick2017overcoming,dhar2019learning}. 3) {Replay-based methods}, which set a fixed-size memory buffer~\cite{buzzega2021rethinking} or generative model~\cite{su2019generative,dedeoglu2023continual,park2020convolutional} to store, produce, and replay historical samples during training. A large number of approaches~\cite{cha2021co2l,chaudhry2021using} have been proposed to improve the performance of replaying. Instead of replaying samples, some approaches~\cite{lao2021two,ho2023prototype} focus on the feature replaying. 4) The prompt-based methods, such as S-prompt~\cite{wang2022s}, L2P~\cite{wang2022learning}, Dualprompt~\cite{wang2022dualprompt}, and CODA-Prompt~\cite{smith2023coda} utilize prompt learning based on a pre-trained model, demonstrating better performance. 5) The optimization-based methods explicitly constrain the updating process of parameters. A typical idea is the gradient projection~\cite{wang2021training,kao2021natural,deng2021flattening} that projects the original gradient in a specified direction for updating parameters more effectively. It is different from the replay-based methods that constrain the gradient by replaying old samples.

In this work, the setup is online continual learning based on the disjoint-task setup with a task-free process. HPCR is proposed as a novel replay-based method to address the forgetting issue. Different from general replay-based methods, HPCR is used to quickly learn novel knowledge and efficiently retain historical knowledge in an online fashion.

% Compared with the general setting, this setup is more practical and difficult.

\subsection{Online Continual Learning}
OCL is a subfield of continual learning that focuses on efficiently learning from a single pass over the online data stream. It is essential in scenarios where a model needs to continuously evolve its knowledge and adapt to new information~\cite{mai2022online}. Researchers in this field focus on developing algorithms and techniques to mitigate catastrophic forgetting, control model adaptation, and efficiently use limited resources to achieve this continuous learning. At present, the replay-based methods are the main solutions of OCL without AOP~\cite{guo2022adaptive}, which is based on orthogonal projection. Although architecture-based methods can be effective for continual learning, they are often impractical for OCL. Since it is unfeasible for large numbers of tasks~\cite{zhu2021class} by the "incremental network architectures". Meanwhile, it is known that regularization-based methods show poor performance in OCL as it is hard to design a reasonable metric to measure the importance of parameters~\cite{zhu2021class}. 

Experience replay (ER)~\cite{rolnick2019experience} that employs reservoir sampling for memory management is usually a strong strategy. Many subsequent methods have been improved based on this method and can generally be divided into three groups. 1) Some approaches belong to \textbf{memory retrieval strategy}, which usually studies how to select more important samples from the buffer for replaying. These representative methods include MIR~\cite{aljundi2019online} using increases in loss, and ASER~\cite{shim2021online} using adversarial shapley value. 2) In the meantime, some approaches~\cite{aljundi2019gradient,jin2021gradient} focus on saving more effective samples to the memory buffer, belonging to the \textbf{memory update strategy}. The core of this type of method is to use the limited samples in the buffer to approximate all previous samples as much as possible. 3) In addition to selecting important samples, it is equally important to conduct efficient anti-forgetting training. The \textbf{model update strategy} is a family of methods ~\cite{mai2021supervised,yin2021mitigating,caccia2021new,guo2022online,gu2022not,chrysakisonline,zhang2022simple,lin2023uer,wei2023online,michel2024learning,kim2024online} to improve the learning efficiency, making the model learn quickly and memory efficiently. All of them belong to the proxy-based replay manners except SCR~\cite{mai2021supervised}, which uses a contrastive-based replay manner. 

The proposed HPCR in this work is a novel model update strategy for OCL. Different from existing approaches, it is a more holistic and effective method proposed based on PCR. After conducting an in-depth analysis of PCR, HPCR improves the model’s feature extraction, generalization, and anti-forgetting capabilities with three components, making PCR better suited for OCL than previous studies.


\subsection{Deep Metric Learning}
Similar to~\cite{yao2022pcl}, our method is inspired by deep metric learning. Deep metric learning aims to measure the similarity among samples by a deep neural network while using an optimal distance metric for learning tasks~\cite{yang2018person}. Generally, the loss functions for this learning task can be categorized into pair-based and proxy-based losses. For one thing, the pair-based methods~\cite{milbich2020sharing,elezi2022group} can mine rich semantic information from anchor-to-sample relations, but converge slowly due to its high training complexity. Since contrastive-based loss is good at learning anchor-to-sample pairs, it can be regarded as a type of pair-based method. For another thing, the proxy-based methods~\cite{zheng2023deep,roth2022non} converge fast and stably, meanwhile may miss partial semantic information by learning anchor-to-proxy relations. The cross-entropy loss can be regarded as a type of proxy-based method.

The proposed HPCR in this work integrates these two types of methods. Different from existing studies, HPCR is proposed to quickly learn novel knowledge and effectively keep historical knowledge for OCL.







